\exercisegroup

\begin{enumerate}
\def\labelenumi{\arabic{enumi})}
\item
  A subset A of a vector space E, is starshaped relative to 0 if for all \(x \in \mathbf{A}\) and every \(\lambda\) such that \(0 \leqslant \lambda < 1\) , the point \(\lambda x\) belongs to A. Let A be starshaped and such that, for each \(x \in \mathbf{A}\) , there exists \(\mu > 1\) such that \(\mu x \in \mathbf{A}\) . Show that if, for every pair of points \(x, y\) of A we have \(\frac{1}{2}(x + y) \in \mathbf{A}\) , then A is convex. Give an example of a non-convex starshaped set A such that \(\frac{1}{2}(\mathbf{A} + \mathbf{A}) \subset \mathbf{A}\) .
\item
  Let A be a convex subset of an affine space E and B a set containing A. Show that, amongst the convex sets that both contain A and are contained in B there exists at least one maximal set; give an example where there are several distinct maximal sets.
\item
  Let A and B be two disjoint convex sets in a vector space E. Show that there exist two disjoint convex sets C, D in E such that \(A \subset C\) , \(B \subset D\) and \(C \cup D = E\) . (Apply th. 2 of S, III, § 2.4 to the set of pairs of disjoint convex sets (M, N) such that \(A \subset M\) and \(B \subset N\) and express the fact that M and N do not meet by the relation \(0 \notin M - N\) . To show that \(C \cup D = E\) , obtain a contradiction supposing that \(x_{0} \notin C \cup D\) ; if \(C'\) (resp. \(D'\) ) is the convex envelope of \(C \cup \{x_{0}\}\) (resp. \(D \cup \{x_{0}\}\) ), show that it is impossible that both \(C' \cap D \neq \varnothing\) and \(C \cap D' \neq \varnothing\) .)
\item
  Let C be a convex cone with vertex 0 in a vector space E; if \((x_{i})_{1\leqslant i\leqslant n}\) is a finite family of points of C such that \(\sum_{i = 1}^{n}\lambda_{i}x_{i} = 0\) for a family of numbers \(\lambda_{i} > 0\) , then C contains the vector subspace of E generated by the \(x_{i}\) .
\item
  Suppose that the vector space E has an enumerably infinite basis \((e_{n})_{n\in\mathbb{N}}\) . Let C be the set of points \(x = \sum_{n} \xi_{n} e_{n}\) such that for the largest index n for which \(\xi_{n} \neq 0\) , we have \(\xi_{n} > 0\) . Show that C is a pointed convex cone such that \(C \cap (-C) = \{0\}\) and \(C \cup (-C) = E\) ; deduce that C is the set of elements of E that are \(\geqslant 0\) for an order structure that is compatible with the vector space structure of E and for which E is linearly ordered. Show that on this ordered vector space the only linear positive form is 0.
\item
  Let E be an affine space of dimension \(\geqslant 2\) and let f be a bijection of E on itself; show that if the image under f of every convex subset of E is also a convex set, then f is an affine linear mapping (consider the inverse mapping of f and note that a closed segment is the intersection of the convex sets which contain its extremities cf.~A, II, § 9, exerc. 7).
\item
  Give an example of two convex sets \(\mathbf{A} \subset \mathbf{R}\) , \(\mathbf{B} \subset \mathbf{R}^2\) , such that the image of the convex set \(\mathbf{A} \times \mathbf{B}\) under the bilinear mapping \((\lambda, x) \mapsto \lambda x\) of \(\mathbf{R} \times \mathbf{R}^2\) in \(\mathbf{R}^2\) is not convex.
\item
  Let \((\mathbf{A}_i)_{1 \leqslant i \leqslant p}\) be a finite family of convex subsets of a vector space \(E\) ; let \(\mathbf{W}_i\) be the subspace obtained by a translation from the affine linear variety generated by \(\mathbf{A}_i\) ( \(1 \leqslant i \leqslant p\) ). If \(\mathbf{W} = \sum_{i=1}^{p} \mathbf{W}_i\) , show that the affine linear variety generated by the convex set \(\sum_{i=1}^{p} \lambda_i \mathbf{A}_i\) (where \(\lambda_i\) are non-zero numbers) is obtained from \(\mathbf{W}\) by a translation.
\end{enumerate}

¶ 9) Let A be a subset of the space \(\mathbf{R}^n\) .

\begin{enumerate}
\def\labelenumi{\alph{enumi})}
\tightlist
\item
  Show that the convex envelope of A is identical with the set of points \(\sum_{i=0}^{n} \lambda_i x_i\) , where \(x_i \in \mathbf{A}\) , \(\lambda_i \geqslant 0\) for \(0 \leqslant i < n\) , and \(\sum_{i=0}^{n} \lambda_i = 1\) . (Establish the following lemma; if \(p + 1\) points \(x_i\) ( \(0 \leqslant i \leqslant p\) ) form an affinely dependent system (that is to say there exists a relation \(\sum_{i=0}^{p} \beta_i x_i = 0\) where the \(\beta_i\) are not all zero and \(\sum_{i=0}^{p} \beta_i = 0\) ) and if \(x = \sum_{i=0}^{p} \alpha_i x_i\) , where the \(\alpha_i\) are \(\geqslant 0\) and \(\sum_{i=0}^{p} \alpha_i = 1\) , then there exists an index \(k \leqslant p\) and \(p\) numbers \(\gamma_i\) ( \(0 \leqslant i \leqslant p\) , \(i \neq k\) ) such that \(\gamma_i \geqslant 0\) for all \(i\) , \(\sum_{i \neq k} \gamma_i = 1\) and \(x = \sum_{i \neq k} \gamma_i x_i\) ; for this compare those of the \(\alpha_i / \beta_i\) that are defined.) b) Let \(a\) be a point of the convex envelope of A which does not belong to the convex envelope of any subset of A with at most \(n\) points. Show then that A contains at least \(n + 1\) connected components. (We can suppose that \(a = 0\) ; let \((b_i)_{0 \leqslant i \leqslant n}\) be a family of \(n + 1\) affinely independent points of A such that 0 belongs to the convex envelope of the \(b_i\) (cf.~a)). For each index \(i\) , let \(C_i\) be the pointed convex cone of vertex 0 generated by the \(b_j\) with indices \(j \neq i\) ; show that A does not meet the frontier of any of the cones - \(C_i\) .) c) If C is a pointed cone with vertex 0 in \(\mathbf{R}^n\) , show that the convex envelope of C is the set of points \(\sum_{i=1}^{n} x_i\) , where \(x_i \in C\) for \(1 \leqslant i \leqslant n\) .
\end{enumerate}

¶ 10) Let C be the convex envelope of a subset A of \(\mathbf{R}^n\) , and let \(a\) be an interior point of C. Show that there exist \(2n\) points \(x_i \in \mathrm{A}\) ( \(1 \leqslant i \leqslant 2n\) ) such that \(a\) is interior to \(\mathrm{C}_0\) , the convex envelope of the \(x_i\) . (Suppose \(a = 0\) , and argue by induction on \(n\) , noting that by exerc. \(9a\) ) there exists a set of \(k + 1\) points \(y_j\) of A ( \(0 \leqslant j \leqslant k\) , \(1 \leqslant k \leqslant n\) ), affinely independent and such that, if V is the affine linear variety generated by the \(y_j\) , then \(0 \in \mathrm{V}\) and, relative to V, 0 is interior to the convex envelope of the set of the \(y_j\) . Then project C on E/V and show that 0 is interior to this projection relative to E/V). Show that in the above statement \(2n\) cannot be replaced by \(2n - 1\) .

\begin{enumerate}
\def\labelenumi{\arabic{enumi})}
\setcounter{enumi}{10}
\tightlist
\item
  \begin{enumerate}
  \def\labelenumii{\alph{enumii})}
  \tightlist
  \item
    Show that, in the space \(\mathbf{R}^n\) , every convex set \(\mathbf{A}\) of dimension \(n\) contains at least one interior point (consider an affinely independent system of \(n + 1\) points of \(\mathbf{A}\) ). Deduce that if \(\mathbf{A}\) is everywhere dense in \(\mathbf{R}^n\) then \(\mathbf{A} = \mathbf{R}^n\) .
  \end{enumerate}
\end{enumerate}

\begin{enumerate}
\def\labelenumi{\alph{enumi})}
\setcounter{enumi}{1}
\item
  Let E be the normed space \(l^1(\mathbf{N})\) of absolutely convergent series of real numbers \(x = (\xi_n)\) (I, p.~4); show that the set P of \(x\) , such that \(\xi_n \geqslant 0\) for every index \(n\) , is a proper convex cone, which generates E but does not contain any interior point.
\item
  Let E be a Hausdorff topological vector space on which there exists a non-continuous linear form \(f\) (cf.~II, p.~86, exerc. 17, a)). Show that the sets A and B defined by the relations \(f(x) \geqslant 0\) , \(f(x) < 0\) are convex, non-empty, complementary, everywhere dense and that each of them generates E (algebraically).
\end{enumerate}

\begin{enumerate}
\def\labelenumi{\arabic{enumi})}
\setcounter{enumi}{11}
\item
  Show that in the space \(R^{n}\) , a necessary and sufficient condition that a convex set should be closed, is that its intersection with every straight line should be closed (cf.~II, p.~74, exerc. 5).
\item
  Show that in the space \(\mathbf{R}^n\) , every non-empty open convex set is homeomorphic to \(\mathbf{R}^n\) (use exerc. 12 of GT, VI, § 2).
\item
  Let A be a non-empty closed convex subset of E, a Hausdorff topological vector space.
\end{enumerate}

\begin{enumerate}
\def\labelenumi{\alph{enumi})}
\item
  Show that, for every \(a \in \mathbf{A}\) , the set \(\bigcap_{\lambda > 0} \lambda(\mathbf{A} - a)\) is a closed convex cone in \(\mathbf{E}\) , with vertex 0, independent of \(a\) . It is called the asymptotic cone of \(A\) and written \(C_A\) . For every \(a \in A\) , the set \(a + C_A\) is the union of \(\{a\}\) and those open half lines that are contained in \(A\) and have \(a\) as an end point.
\item
  If \(x, y\) are two points of \(\mathbf{A}\) such that \((x + \mathrm{C}_{\mathbf{A}}) \cap (y + \mathrm{C}_{\mathbf{A}})\) is a cone whose vertex \(z \in \mathbf{A}\) , then this cone is necessarily \(z + \mathrm{C}_{\mathbf{A}}\) .
\item
  If B is a second closed convex subset of E such that \(A \cap B \neq \varnothing\) , then \(C_{A \cap B} = C_{A} \cap C_{B}\) . d) Let \(V_{A}\) be the largest vector subspace (necessarily closed in E) which is contained in \(C_{A}\) . Show that if \(\phi\) is the canonical homomorphism of E on \(E/V_{A}\) , then \(A = \phi^{-1}(A_{0})\) , where \(A_{0}\) is a closed convex set in \(E/V_{A}\) which does not contain any straight line.
\item
  In the Banach space \(\mathcal{B}(\mathbf{N})\) of bounded mappings of N in R (I, p.~4) give an example of a closed convex non-bounded set A, for which \(C_{A} = \{0\}\) and which is such that for every \(b \neq 0\) in E, there exists an integer k for which \((\mathbf{A} + kb) \cap \mathbf{A} = \varnothing\) .
\end{enumerate}

\begin{enumerate}
\def\labelenumi{\arabic{enumi})}
\setcounter{enumi}{14}
\tightlist
\item
  \begin{enumerate}
  \def\labelenumii{\alph{enumii})}
  \tightlist
  \item
    Let A be a closed convex subset of a Hausdorff topological vector space E. If for some point \(x_0 \in \mathbf{A}\) there exists a neighbourhood \(\mathbf{V}\) of \(x_0\) in \(\mathbf{E}\) such that \(\mathbf{V} \cap \mathbf{A}\) is compact, then show that \(\mathbf{A}\) is locally compact. Deduce that the closure in \(\mathbf{E}\) of a locally compact convex set is locally compact.
  \end{enumerate}
\end{enumerate}

\begin{enumerate}
\def\labelenumi{\alph{enumi})}
\setcounter{enumi}{1}
\tightlist
\item
  Let A be a closed convex set that is locally compact but not compact in E; show that the asymptotic cone (exerc. 14) is not the single point \(\{0\}\) .
\end{enumerate}

¶ 16) Let A, B be two closed convex subsets of a Hausdorff topological vector space E. Suppose further that B is locally compact and that \(C_A \cap C_B = \{0\}\) . Show that A - B is closed in E. (Let \(b \in B\) , and W be a closed neighbourhood of 0 in E such that \(B \cap (b + W)\) is compact. Let \(c \in \overline{A - B}\) ; for every neighbourhood V of 0 in E consider the set \(M_V\) of those \(y \in B\) such that \(A \cap (c + y + V) \neq \emptyset\) . Consider two cases according as to whether there exists a V for which \(M_V\) is relatively compact, or there does not exist such a V; in the second case, consider the filter base formed from the sets \(P_{V,n} = M_V \cap \complement (b + nW)\) where V varies in the set of closed neighbourhood of 0 in E and n varies in N; form the cone with vertex b generated by \(P_{V,n}\) and its intersection with the frontier of \(b + W\) ).

¶ 17) In a Hausdorff topological vector space E, a closed convex set A is said to be parabolic if, for every \(z \notin A\) , each half-line originating at \(z\) and contained in \(z + C_A\) meets A.

\begin{enumerate}
\def\labelenumi{\alph{enumi})}
\tightlist
\item
  Give an example of a parabolic convex set A in \(\mathbf{R}^2\) such that \(C_A\) is not just a single half-line. b) Let A be a closed convex set in E such that \(C_A \neq \{0\}\) , but such that A is not parabolic. Show that if \(z \notin A\) is such that \(z + C_A\) contains a half-line D with end point \(z\) which does not meet A, then neither the convex envelope of \(A \cup \{z\}\) nor the pointed cone with vertex \(z\) generated by A is closed in E.
\end{enumerate}

Further if \(\mathbf{D}'\) is the closed half-line originating at \(z\) and opposite to \(\mathbf{D}\) (so that \(\mathbf{D} = 2z - \mathbf{D}'\) ) then \(\mathbf{D}' + \mathbf{A}\) is not closed in \(\mathbf{E}\) , and there exists a plane \(\mathbf{P}\) containing \(\mathbf{D}\) and a closed convex set \(\mathbf{B} \subset \mathbf{P}\) , such that \(\mathbf{B} \cap \mathbf{A} = \varnothing\) but that the distance of \(\mathbf{B}\) from \(\mathbf{P} \cap \mathbf{A}\) (in any norm on \(\mathbf{P}\) ) is zero.

\begin{enumerate}
\def\labelenumi{\alph{enumi})}
\setcounter{enumi}{2}
\item
  In E, let A be a closed convex set that is locally compact and parabolic; show that if \(\mathbf{B} \subset \mathbf{E}\) , is closed and convex then \(\mathbf{A} - \mathbf{B}\) is closed (same method as exerc. 16).
\item
  Let A, \(A'\) be two closed convex subsets of E that are locally compact and parabolic; show that the convex envelope of \(A \cup A'\) is closed in E (same method as in exerc. 16). Give an example in \(\mathbf{R}^2\) where \(\mathbf{A}\) is parabolic, \(\mathbf{A}'\) is non-parabolic and the convex envelope of \(\mathbf{A} \cup \mathbf{A}'\) is not closed in \(\mathbf{R}^2\) .
\item
  In E, let A be a closed convex set that is locally compact and parabolic; show that for all \(z \notin \mathbf{A}\) , the pointed convex cone with vertex \(z\) , generated by A, is closed in E (use \(d\) ).
\end{enumerate}

\(f\) ) Let \(\mathbf{E}_1, \mathbf{E}_2\) be two Hausdorff topological spaces, \(\mathbf{A}_1\) (resp. \(\mathbf{A}_2\) ) a closed convex set in \(\mathbf{E}_1\) (resp. \(\mathbf{E}_2\) ) that is also parabolic. Show that the set \(\mathbf{A}_1 \times \mathbf{A}_2\) is parabolic in \(\mathbf{E}_1 \times \mathbf{E}_2\) .

* g) Show that a barrelled space of infinite dimension does not contain a parabolic closed convex set that is both locally compact and not compact.

\(h\) ) Let \(\mathrm{E}_0 = l^2 (\mathbf{N})\) , and in \(\dot{\mathrm{E}}_0\) , let \(\mathbf{K}\) be the set of points \(x = (\xi_n)\) such that \(|\xi_n| \leqslant 1 / (n + 1)\) for all \(n\) ; \(\mathbf{K}\) is compact. The cone E of vertex 0 generated by \(\mathbf{K}\) is a vector subspace of \(\mathrm{E}_0\) and \(\mathbf{K}\) is absorbent in E. Let \(p\) be the gauge of \(\mathbf{K}\) in E; it is a lower semi-continuous function. In the normed product space \(\mathrm{E} \times \mathbf{R}\) , show that the set A of points \((x, \zeta)\) such that \(\zeta \geqslant (p(x))^2\) is closed, convex, parabolic and locally compact. Show that \(\mathrm{A} + (-\mathrm{A})\) and the convex envelope of \(\mathrm{A} \cup (-\mathrm{A})\) are not locally compact. *

\begin{enumerate}
\def\labelenumi{\arabic{enumi})}
\setcounter{enumi}{17}
\item
  Let \(\mathbf{C}\) be a proper closed convex cone of vertex 0 in \(\mathbf{R}^n\) . Show that the complement of the set \(\mathbf{C} \cap \mathbf{S}_{n-1}\) on the sphere \(\mathbf{S}_{n-1}\) is homeomorphic to \(\mathbf{R}^{n-1}\) (make a stereographic projection from a point of \(\mathbf{C} \cap \mathbf{S}_{n-1}\) , and use exerc. 12 of GT, VI, § 2). If \(\mathbf{C}\) contains an interior point, show that \(\mathbf{C} \cap \mathbf{S}_{n-1}\) , is homeomorphic to the closed ball \(\mathbf{B}_{n-1}\) (same method).
\item
  \begin{enumerate}
  \def\labelenumii{\alph{enumii})}
  \tightlist
  \item
    Let \(\mathbf{A}\) be an unbounded closed convex set in \(\mathbf{R}^n\) , that does not contain any line, but does contain an interior point. Show that the frontier of \(\mathbf{A}\) is homeomorphic to \(\mathbf{R}^{n-1}\) (use exercs. 15, b) and 18).
  \end{enumerate}
\end{enumerate}

\begin{enumerate}
\def\labelenumi{\alph{enumi})}
\setcounter{enumi}{1}
\tightlist
\item
  In a Hausdorff topological vector space E, let A be a closed convex set that does not contain any line and is of dimension \(\geqslant 2\) . Show that the frontier of A is connected (use a) and GT, VI, § 2, exerc. 12).
\end{enumerate}

\begin{enumerate}
\def\labelenumi{\arabic{enumi})}
\setcounter{enumi}{19}
\tightlist
\item
  \begin{enumerate}
  \def\labelenumii{\alph{enumii})}
  \tightlist
  \item
    In a vector space E, let A be a convex set that generates E and meets every straight line in a set that is closed relative to the straight line. Show that the following conditions are equivalent:
  \end{enumerate}
\end{enumerate}

\(\alpha)\) There exists a line D such that D meets A in a compact segment that is not empty.

β) There exists a line D such that every line parallel to D meets A in a compact segment.

\(\gamma)\) A is distinct from E and is not a half-space determined by a hyperplane of \(\tilde{\mathbf{E}}\) .

(To show that \((\gamma) \Rightarrow (\alpha)\) use exerc. 14, \(d\) ) of II, p.~67, and reduce to the case \(\mathbf{E} = \mathbf{R}^2\) .)

\begin{enumerate}
\def\labelenumi{\alph{enumi})}
\setcounter{enumi}{1}
\tightlist
\item
  In a Hausdorff topological vector space E, let A be a closed convex set which contains an interior point. Show that if the frontier of A is a non-empty linear variety, then A is a closed half-space (use exerc. 14, d) of II, p.~67, to show that the frontier of A is necessarily a hyperplane, then apply a)).
\end{enumerate}

¶ 21) a) Let \(\mathrm{A}_i\) ( \(1 \leqslant i \leqslant r\) ) \(r > n + 1\) , be a family of convex subsets of \(\mathbf{R}^n\) such that any \(r - 1\) of the \(\mathrm{A}_i\) have a non-empty intersection; show that the \(r\) sets \(\mathrm{A}_i\) have a non-empty intersection (Helly's theorem). (Let \(x_i\) be a point of the intersection of the \(\mathrm{A}_j\) with indexes \(j \neq i\) ;

there exist \(r\) numbers \(\lambda_{i}\) which are not all zero and are such that \(\sum_{i=1}^{r} \lambda_{i} = 0\) and \(\sum_{i=1}^{r} \lambda_{i} x_{i} = 0\) ;

in this last equation take to one side those terms with \(\lambda_{i} \geqslant 0\) and to the other those with \(\lambda_{i} < 0\) .)

\begin{enumerate}
\def\labelenumi{\alph{enumi})}
\setcounter{enumi}{1}
\item
  Given a family of compact convex sets in \(\mathbf{R}^n\) , show that the intersection of all the sets of the family is non-empty if the intersection of any selection of \(n + 1\) sets of the family is non-empty.
\item
  In \(\mathbf{R}^n\) , let \(\mathbf{K}\) be a convex set and \((\mathbf{A}_i)_{1 \leqslant i \leqslant r}\) be a family of \(r > n + 1\) convex sets. Suppose that for every selection of \(n + 1\) indices \((i_k)\) each less than or equal to \(r\) there exists \(a \in \mathbf{R}^n\) such that \(a + \mathbf{K}\) contains each of the \(A_{i_k}\) . Show that then there exists \(b \in \mathbf{R}^n\) such that \(b + \mathbf{K}\) contains all the \(\mathbf{A}_i\) . Show that similar results hold if « contains » is replaced by « is contained in » or by « meets in a non-empty set ». (For each index \(i\) , consider the set \(\mathbf{C}_i\) of the \(x \in \mathbf{R}^n\) for which \(x + \mathbf{K} \supset \mathbf{A}_i\) (or \(x + \mathbf{K} \subset \mathbf{A}_i\) , or \((x + \mathbf{K}) \cap \mathbf{A}_i \neq \emptyset\) )). Generalize to any family of compact convex sets of \(\mathbf{R}^n\) .
\end{enumerate}

\begin{enumerate}
\def\labelenumi{\arabic{enumi})}
\setcounter{enumi}{21}
\item
  In \(\mathbf{R}^2\) consider a set of \(2m\) points of the form \((a_i, b_i')\) , \((a_i, b_i'')\) where \(b_i' \leqslant b_i''\) for \(1 \leqslant i \leqslant m\) . Let \(n\) be an integer \(< m - 2\) . In order that there should exist a polynomial \(\mathrm{P}(x)\) of degree \(\leqslant n\) such that \(b_i' \leqslant \mathrm{P}(a_i) \leqslant b_i''\) for \(1 \leqslant i \leqslant m\) , it is sufficient that, for every family \((i_k)_{1 \leqslant k \leqslant n + 2}\) of \(n + 2\) indices \(i\) , there exists a polynomial \(\mathrm{Q}(x)\) of degree \(\leqslant n\) such that \(b_{i_k}' \leqslant \mathrm{Q}(a_{i_k}) \leqslant b_{i_k}''\) for every integer \(k\) such that \(1 \leqslant k \leqslant n + 2\) . (Use exerc. 21, a).)
\item
  Show that in a topological vector space, the convex envelope of an open set is an open set.
\item
  Let \(\mathbf{M}\) be an everywhere dense convex set in a topological vector space \(\mathrm{E}\) (cf.~II, p.~66, exerc. 11, c)); show that, for every closed hyperplane \(\mathbf{H}\) in \(\mathbf{E}\) , the set \(\mathbf{H} \cap \mathbf{M}\) is dense in \(\mathbf{H}\) (for every point \(x_0 \in H\) , and every balanced neighbourhood \(\mathbf{V}\) of 0 in \(\mathbf{E}\) , consider the intersections of \(x_0 + \mathbf{V}\) and the two open half-spaces determined by \(\mathbf{H}\) , and deduce that \(x_0 + \mathbf{V} + \mathbf{V}\) meets \(\mathbf{H} \cap \mathbf{M}\) ).
\item
  \begin{enumerate}
  \def\labelenumii{\alph{enumii})}
  \tightlist
  \item
    Show that, in a topological vector space, every convex set with an interior point, is such that its frontier is nowhere dense (use prop. 16 of II, p.~14).
  \end{enumerate}
\end{enumerate}

\begin{enumerate}
\def\labelenumi{\alph{enumi})}
\setcounter{enumi}{1}
\tightlist
\item
  In a Hausdorff topological vector space E, let A be a closed convex set with an interior point, and let H be a closed hyperplane that contains an interior point of A. Show that the intersection of H and of the frontier F of A is a set which is nowhere dense relative to F (to show that in every neighbourhood of a point of \(\mathbf{H} \cap \mathbf{F}\) there exist points of F not in H, reduce to the case when E is of dimension 2).
\end{enumerate}

¶ 26) In a Hausdorff topological vector space E, let A be a connected closed set with the following property : for every \(x \in A\) , there exists a closed neighbourhood V of \(x\) in E such that V ∩ A is convex. Show that A is convex. For this establish the following statements.

\begin{enumerate}
\def\labelenumi{\alph{enumi})}
\item
  Show that any two points of A can be joined by a broken line in A (same method as GT, VI, § 1, exerc. 6).
\item
  Show that, if two points in A can be joined by a broken line in A with n > 1 segments, then they can also be joined in A by a broken line with n - 1 segments. (Induction on n reduces to the case n = 2 which is equivalent to taking \(R^{2}\) as E; then let T be a triangle with vertices a, b, c such that the closed segments ac, bc are contained in A, but the closed segment ab is not; consider a point of the closure of the intersection of \(\complement A\) and the interior of T that is farthest from the line ab, and show that the existence of such a point contradicts the hypothesis.)
\end{enumerate}

¶ 27) a) Let B be a non-empty closed convex subset of E, a Hausdorff topological vector space, and let X be a non-empty compact set in E. Show that if A is a subset of E such that \(A + X \subset B + X\) , then \(A \subset B\) (if \(a \in A\) , consider a sequence \((x_{n})\) of points of X defined inductively by the relation \(a + x_{n} = b_{n} + x_{n+1}\) , where \(b_{n} \in B\) ). Deduce that, if A, B are two non-empty subsets of E, using the distance in E and the procedure of GT, IX, § 2, exerc. 6. Show \(A + X = B + X\) implies the relation A = B.

\begin{enumerate}
\def\labelenumi{\alph{enumi})}
\setcounter{enumi}{1}
\item
  Let E be a normed space, \(\sigma\) the distance function defined on the set \(\mathfrak{F}(E)\) of closed non-empty subsets of E, using the distance in E and the procedure of TG, IX, p.~91, exerc. 6. Show that if A, B, C are three non-empty compact convex sets in E then \(\sigma(A+C, B+C)=\sigma(A, B)\) (if \(S_{\lambda}\) is the ball defined by \(\|x\|\leqslant\lambda\) , note that \(A+S_{\lambda}\) and \(B+S_{\lambda}\) are closed convex sets and use a)).
\item
  Deduce from a) and b) that the set \(\Re(\mathrm{E})\) of non-empty, compact, convex subsets of a normed space E, with the distance \(\sigma\) , can be identified with a cone in a normed space of which the laws of composition induce on \(\Re(\mathrm{E})\) the laws \((\mathrm{A}, \mathrm{B}) \to \mathrm{A} + \mathrm{B}\) and \((\lambda, \mathrm{A}) \to \lambda\mathrm{A}\) .
\end{enumerate}

\begin{enumerate}
\def\labelenumi{\arabic{enumi})}
\setcounter{enumi}{27}
\tightlist
\item
  Let \(f\) be a convex function defined over the convex subset \(A\) of a vector space \(E\) .
\end{enumerate}

\begin{enumerate}
\def\labelenumi{\alph{enumi})}
\item
  Show that if \(\mathbf{A}\) is absorbent and \(f\) is non-constant then \(f\) cannot attain its upper bound in \(\mathbf{A}\) at the point 0.
\item
  Show that the subset of points of A, at which \(f\) attains its lower bound in A, is convex.
\end{enumerate}

¶ 29) Let E be a Hausdorff topological vector space, and C be a non-empty open convex non-pointed cone with vertex 0, in E. A convex neighbourhood of 0 in E is denoted by V. If f is a convex function that is defined and bounded above in C ∩ V, show that \(f(x)\) tends to a finite limit as x tends to 0, where \(x \in C \cap V\) . (Let \(\beta = \limsup_{x \to 0, x \in C \cap V} f(x)\) ; obtain a contradiction,

supposing that for some \(\alpha > 0\) , and every neighbourhood \(\mathbf{W}\) of \(0\) , there exists a point \(y \in \mathbf{C} \cap \mathbf{V} \cap \mathbf{W}\) such that \(f(y) < \beta - \alpha\) . Show that there exists \(a \in \mathbf{C} \cap \mathbf{V}\) such that \(f(\rho a) \geqslant \beta - \frac{1}{2}\alpha\) for \(0 < \rho \leqslant 1\) ; deduce that, on a line joining a point of the form \(\rho a\) ( \(\rho\) sufficiently small) to a point \(y\) in \(\mathbf{C} \cap \mathbf{V}\) such that \(f(y) < \beta - \alpha\) , and that is sufficiently close to \(0\) , there exist points of \(\mathbf{C} \cap \mathbf{V}\) where \(f\) is arbitrarily large.)

\begin{enumerate}
\def\labelenumi{\arabic{enumi})}
\setcounter{enumi}{29}
\tightlist
\item
  \begin{enumerate}
  \def\labelenumii{\alph{enumii})}
  \tightlist
  \item
    Give an example of a convex function that is defined over a compact convex subset K of \(R^{2}\) , that is bounded and lower semi-continuous in K, but is not continuous at a point of the frontier of K (consider the gauge of a disc of which 0 is a frontier point).
  \end{enumerate}
\end{enumerate}

\begin{enumerate}
\def\labelenumi{\alph{enumi})}
\setcounter{enumi}{1}
\item
  Deduce from \(a\) an example of a convex function defined in an open half-plane \(\mathbf{D}\) of \(\mathbb{R}^2\) , not bounded above in \(\mathbf{D}\) and not tending to a limit at a frontier point of \(\mathbf{D}\) .
\item
  Deduce from a) an example of a convex lower semi-continuous function defined over a compact convex subset A of \(R^{2}\) but not bounded above in A. (Take for A the set of points \((\xi, \eta)\) such that \(\xi^{4} \leqslant \eta \leqslant 1\) in \(R^{2}\) .)
\end{enumerate}

¶ 31) Let \(x_0\) be a point in the closure of A, a non-empty convex subset of a Hausdorff topological vector space E. Let \(f\) be a convex function defined over A. Use \(\mathfrak{D}\) to denote the set of closed half lines D which originate at \(x_0\) , for which \(A \cap D\) contains an open segment with end point \(x_0\) . The union C of the half lines \(D \in \mathfrak{D}\) is a convex cone with vertex \(x_0\) .

\begin{enumerate}
\def\labelenumi{\alph{enumi})}
\item
  Show that, for each fixed \(\mathbf{D} \in \mathfrak{D}\) , as \(x\) tends to \(x_0\) such that \(x \in \mathbf{D} \cap \mathbf{A}\) and \(x \neq x_0\) , either \(f(x)\) tends to a finite limit or to \(+\infty\) .
\item
  Let \(\mathfrak{I}\) be the subset of those \(D \in \mathfrak{D}\) , for which the limit of \(f(x)\) in \(a\) is \(+\infty\) ; if \(x_0 \in A\) then \(\mathfrak{I}\) is empty. Show that \(\mathfrak{I}\) cannot contain two opposite half lines; if \(D\) and \(D'\) are two distinct half lines in \(\mathfrak{I}\) and \(P\) is the plane determined by \(D\) and \(D'\) then, either, every half line \(D''\) of \(\mathfrak{D}\) in \(P\) belongs to \(\mathfrak{I}\) , or \(D\) and \(D'\) are the only two half lines of \(\mathfrak{I}\) lying in \(P\) . Deduce that if \(\mathfrak{I} \neq \mathfrak{D}\) , then no half line \(D \in \mathfrak{I}\) contains an internal point (II, p.~26) of the cone \(C\) relative to the vector subspace generated by \(C\) .
\item
  Let \(\mathfrak{F}\) be the set of half lines in \(\mathfrak{D}\) that are not in \(\mathfrak{I}\) . Show that the union of the half lines of \(\mathfrak{F}\) is a convex cone, and for each half line \(\mathrm{D} \in \mathfrak{F}\) the limit of \(f(x)\) defined in \(a\) ) is independent of \(\mathrm{D}\) (use exerc. 29 above); further if \(x_0 \in \mathrm{A}\) this limit is \(\leqslant f(x_0)\) , and it is equal to \(f(x_0)\) when \(\mathfrak{F}\) contains two opposite half lines.
\end{enumerate}

\(d\) ) Let \(f\) be a non-continuous linear form over \(\mathbf{E}\) (cf.~II, p.~86, exerc. 17, \(a\) ) and take \(\mathbf{A} = \mathbf{E}\) ; show that every closed half line, originating at \(x_0\) , belongs to \(\mathfrak{F}\) , but that

\[
\lim _ {x \to x _ {0}} \inf f (x) = - \infty \text { and } \lim _ {x \to x _ {0}} \sup f (x) = + \infty
\]

(use prop. 21 of II, p.~18).

\begin{enumerate}
\def\labelenumi{\arabic{enumi})}
\setcounter{enumi}{31}
\item
  Let \(\mathbf{K}\) be a compact convex set in a Hausdorff topological vector space \(\mathbf{E}\) and let \(f\) be an upper semi-continuous convex function defined over \(\mathbf{K}\) . Show that \(f\) is bounded over \(\mathbf{K}\) . (Observe first that \(f\) is bounded above in \(\mathbf{K}\) ; if \(f\) is not bounded below show that \(\liminf_{y\to x,y\neq x}f(y) = -\infty\) for every point \(x\in \mathbf{K}\) , and that this contradicts Baire's theorem.) Give an example where \(f\) is not continuous.
\item
  Let E be a finite dimensional Hausdorff topological vector space, and let K be a compact convex subset of E. Show that every convex function defined over K is bounded below in K (compare exerc. 31, d)).
\item
  Let U, V be two open convex sets in a Hausdorff topological vector space E such that \(\overline{V} \subset U\) and that U does not contain any half line. Let F be a set of convex functions defined in \(\overline{U}\) , uniformly bounded above on the frontier of U and uniformly bounded below on the frontier of V. Show that F is equicontinuous.
\item
  Let U be a non-empty open convex set in \(R^{n}\) and F be a set of convex functions defined over U. Let \(\Phi\) be a filter on F that converges pointwise in U to a finite function \(f_{0}\) ; show that \(\Phi\) converges uniformly to \(f_{0}\) in every compact subset of U (use exerc. 34).
\item
  Let A be a compact convex set in \(\mathbf{R}^n\) and B its projection on the subspace \(\mathbf{R}^{n-1}\) (identified as the hyperplane with equation \(\xi_n = 0\) ). Show that there exist two convex functions \(f_1, f_2\) defined over B, such that A is identical with the set of points \((x, \zeta)\) of \(\mathbf{R}^n\) where \(x \in \mathbf{B}, y \in \mathbf{R}\) and \(f_1(x) \leqslant \zeta \leqslant -f_2(x)\) .
\item
  Let E be a vector space; in order that a convex set F of \(E \times R\) should be formed of pairs \((x, \zeta)\) such that \(f(x) \leqslant \zeta\) (resp. \(f(x) < \zeta\) ) for a convex function f defined over a convex subset X of E, it is necessary and sufficient that the projection of F on E should be identical with X and that, for all \(x \in X\) , the set \(F(x)\) of F that projects onto x should be a closed (resp. open) interval unlimited to the right (i.e.~not bounded above).
\item
  Let \(\mathbf{X}\) be a convex set of an affine space \(\mathbf{E}\) and \(p\) an affine linear mapping of \(\mathbf{E}\) in a second affine linear space \(\mathbf{E}_1\) . Write \(\mathbf{X}_1 = p(\mathbf{X})\) . For every real-valued function \(f\) defined in \(\mathbf{X}\) and every \(x_1 \in \mathbf{X}_1\) , let
\end{enumerate}

\[
f _ {1} (x _ {1}) = \inf _ {p (x) = x _ {1}} f (x).
\]

Show that if \(f\) is convex and if \(f_{1}(x_{1}) > -\infty\) for all \(x_{1} \in X_{1}\) , then \(f_{1}\) is a convex function.

\begin{enumerate}
\def\labelenumi{\arabic{enumi})}
\setcounter{enumi}{38}
\tightlist
\item
  Let \(\mathbf{E}\) be a finite dimensional Hausdorff topological vector space.
\end{enumerate}

\begin{enumerate}
\def\labelenumi{\alph{enumi})}
\item
  Let \(\mathfrak{F}(\mathrm{E})\) be the family of closed non-empty sets of E, carrying the uniform structure deduced from the uniform structure of E by the procedure of GT, II, § 1, exerc. 5, a). Show that, the set \(\mathfrak{C}(\mathrm{E})\) of non-empty closed convex sets of E, is closed in the space \(\mathfrak{F}(\mathrm{E})\) . Deduce that if K is a compact set in E, the set of non-empty closed convex sets in E that are contained in K, is a compact set in \(\mathfrak{C}(\mathrm{E})\) (cf.~GT, § 4, exerc. 11).
\item
  Let \(\mathfrak{R}_0(\mathrm{E})\) be the set of compact convex subsets of \(\mathrm{E}\) that contain 0 as an interior point. For every set \(\mathrm{A} \in \mathfrak{R}_0(\mathrm{E})\) , let \(p_{\mathrm{A}}\) be the gauge of \(\mathrm{A}\) (II, p.~20). Show that \(\mathrm{A} \mapsto p_{\mathrm{A}}\) is an isomorphism of the uniform subspace \(\mathfrak{R}_0(\mathrm{E})\) of \(\mathfrak{C}(\mathrm{E})\) on a subspace of the space \(\mathcal{C}_{c}(\mathrm{E};\mathbf{R})\) of continuous real valued functions in \(\mathrm{E}\) , carrying the uniform structure of compact convergence (GT, X, § 1.6).
\end{enumerate}

\begin{enumerate}
\def\labelenumi{\arabic{enumi})}
\setcounter{enumi}{39}
\item
  In a topological vector space E, let U be a convex neighbourhood of \(x_{0}\) and let f be a real-valued continuous convex function in U. Show that there exists a convex neighbourhood \(V \subset U\) of \(x_{0}\) and a convex continuous function \(f_{1}\) in E such that \(f_{1}|V = f|V\) .
\item
  Let H be a hyperplane in a vector space E that does not contain 0 and let S be a convex set contained in H.
\end{enumerate}

\begin{enumerate}
\def\labelenumi{\alph{enumi})}
\item
  Suppose that the intersection of S with each line in H is a compact segment. Let \(a, b\) be two distinct points of E such that there exist two numbers \(\lambda > 0\) , \(\mu > 0\) for which \(b + \mu S \subset a + \lambda S\) ; show that if \(c\) is the point where the line joining \(a\) and \(b\) meets the hyperplane H' parallel to H which contains \(a + \lambda S\) , then \(c \in b + \mu S\) and \(b + \mu S\) is the image of \(a + \lambda S\) by a homothety of centre \(c\) transforming \(a\) into \(b\) . (Reduce to the case where E is of dimension 2.)
\item
  With the same hypotheses on S, let a, b be two distinct points of E and suppose that there exists a point \(c \in E\) and three numbers \(\lambda > 0\) , \(\mu > 0\) , \(v \geqslant 0\) such that \((a + \lambda S) \cap (b + \mu S) = c + vS\) . Show that if A (resp. B) is the cone with vertex a (resp. b) generated by \(a + \lambda S\) (resp. \(b + \mu S\) ), and \(H''\) the hyperplane parallel to H passing through c, then \(H'' \cap A \cap B = \{c\}\) (use a)). c) Suppose that H is the affine linear variety generated by S. Let C be the cone with vertex 0 generated by S. Show that the following two conditions are equivalent:
\end{enumerate}

α) E is a lattice for the order on E of which C is the set of elements ≥ 0.

β) For any points \(x, y\) of \(E\) and numbers \(\lambda > 0\) , \(\mu > 0\) such that the set \((x + \lambda S) \cap (y + \mu S)\) is not empty, there exists \(z \in E\) and \(\nu \geqslant 0\) such that this set is \(z + \nu S\) .

(To prove that \(\alpha\) ) implies \(\beta\) ), reduce to the case \(y = 0\) and use the fact that if \((s_i)\) is a finite family of points of S and \((\lambda_i)\) is a family of real numbers such that \(\sum_{i} \lambda_i s_i = 0\) , then \(\sum_{i} \lambda_i = 0\) .

To prove that \(\beta)\) implies \(\alpha)\) use \(b)\) .

When S satisfies the equivalent conditions \(\alpha)\) and \(\beta)\) , we say that S is a simplex in E. When E is of finite dimension, the convex envelope of a finite set of points affinely independent in H and generating H is a simplex * (the converse is also true : cf.~INT, II, 2nd ed., § 2, exerc. 7)).

\begin{enumerate}
\def\labelenumi{\arabic{enumi})}
\setcounter{enumi}{41}
\item
  Generalise the prop. 18 of II, p.~16 to the case of an ordered set E with a Hausdorff topology for which the intervals \(\{a, \to (\text{and}) \gets, a\}\) are closed for all \(a \in E\) .
\item
  Let K be a complete valued division ring of which the absolute value is an ultrametric. In a left vector space E on K, we say that a set A is ultraconvex if the relations \(x \in A\) , \(y \in A\) , \(|\lambda| \leqslant 1\) , \(|\mu| \leqslant 1\) imply \(\lambda x + \mu y \in A\) .
\end{enumerate}

\begin{enumerate}
\def\labelenumi{\alph{enumi})}
\tightlist
\item
  Generalize the prop. 1, 2, 5, 6, 7 of II, p.~8 and p.~9. Show that the smallest ultraconvex set containing a given set \(\mathbf{M}\) is the set of linear combinations \(\sum_{\iota} \lambda_{\iota} x_{\iota}\) , where \(x_{\iota} \in \mathbf{M}\) and \(|\lambda_{\iota}| \leqslant 1\)
\end{enumerate}

for all \(\iota\) .

\begin{enumerate}
\def\labelenumi{\alph{enumi})}
\setcounter{enumi}{1}
\item
  Suppose that E is a topological vector space over K. Show that the closure of an ultraconvex set is ultraconvex, and that an ultraconvex set with a non-empty interior is open.
\item
  Let A be an absorbent and ultraconvex set in E. Show that if, for all \(x \in E\) , we put \(p(x) = \inf_{x \in \mathbb{P}^{\mathbf{A}}} |\rho|\) , then p is an ultra-semi-norm on E (II, p.~2). Generalize prop. 23 of II, p.~20,
\end{enumerate}

to the case where the absolute value of K is obtained from a discrete valuation.

